To compare the spatial predictions on the same anatomy, we compute all maps
in one open-access subject,
IXI002 of the IXI dataset\footnote{\url{https://brain-development.org/ixi-dataset/}}
(T1-weighted MRI; female, 35 years). The T1 image was segmented into six tissues
(scalp, skull, cerebrospinal fluid, gray matter, white matter and eyes) and meshed with
SimNIBS~3.2.3 \texttt{headreco} \cite{nielsen2018}, with its CAT12 option
\cite{gaser2024}. Tissue conductivities were 0.33 (scalp), 0.008 (skull), 1.79
(cerebrospinal fluid), 0.40 (gray matter), 0.15 (white matter) and 0.33~S/m (eyes)
\cite{salvador2025}. These are low-frequency values; no dispersion correction was
applied at the kHz carriers. The volume conductor is linear and quasi-static, so the
carrier frequencies do not enter the maps.

The electrodes follow the published hippocampal montage of Violante et al.\
\cite{violante2023}, placed by its recipe on the subject's registered 10--10
positions. Pair~1 (anterior) joins e1, 25~mm from T9 toward FT9, to e2, 10~mm from AF8
toward Fp2. Pair~2 (posterior) joins e3, 25~mm from T9 toward TP9, to e4, 10~mm from P8
toward PO8. Each electrode is a 20-mm disc of 3~mm of gel (4.0~S/m) under 2~mm of metal
(10~S/m), meshed into the head model with SimNIBS~3.2.3; the resulting mesh has
$3.64\times10^6$ linear tetrahedra. As a montage sensitivity, the 10--10 rendering of
this montage, FT7--F8 (anterior) and TP7--P8 (posterior), which Yatsuda et al.\
\cite{yatsuda2024} used as their starting montage, was meshed and solved in the same
way. Each montage has its own mesh, carrying only its own four electrodes.

Each pair was solved once at 1~A. The potential $\phi_j$ (V) satisfies
$\nabla\!\cdot\!(\sigma_{\rm cond}\nabla\phi_j)=0$, where
$\sigma_{\rm cond}$ is tissue conductivity (S/m). Current enters and leaves through each
electrode's metal contact as a Neumann condition, distributed in proportion to nodal
area; the rest of the scalp carries no normal current, and one pinned node fixes the
gauge. The linear finite-element system was assembled and solved with Neuroelectrics'
temporal-interference library, by conjugate gradients with algebraic-multigrid
preconditioning, to a relative residual of $10^{-8}$. The field
$-\nabla\phi_j$ was recovered at gray- and white-matter
nodes by superconvergent patch recovery \cite{zienkiewicz1992}, separately within each
compartment (446{,}877 node--compartment pairs). Dividing these fields by the solve
current gives the field-per-current vectors $\mathbf F_j$, expressed in
(V/m)/mA to match $I_j^{\rm pair}$. Every current condition is a scaling of these two fields.
With peak pair currents $I_1^{\rm pair}$ (anterior) and
$I_2^{\rm pair}$ (posterior) summing to 4~mA, the ratio 1:1 is $2+2$~mA and 1:3 is
1~mA anterior plus 3~mA posterior, the two conditions of Violante et al. The ratio sweep
sets the posterior share $I_2^{\rm pair}/(I_1^{\rm pair}+I_2^{\rm pair})$ from 0.1 to
0.9 in steps of 0.025.

The field-to-polarization conversion follows \S\ref{sec:envelope}:
$\varepsilon_j(\mathbf x)=I_j^{\rm pair}\mathcal P_j(\mathbf x)$, with
$\mathcal P_j=\lambda_0\,\mathbf F_j\!\cdot\!\mathbf n$, one uniform
static coupling $\lambda_0=0.5$~mV per V/m \cite{radman2009}, and
$\lambda_f\approx\lambda_0$. A uniform coupling multiplies the envelope map by
$\lambda_0$ and the two second-order maps by $\lambda_0^2$. Target-to-reference ratios,
normalized maps and focus positions do not depend on it. Hippocampal couplings of
$0.12$--$0.18$~mV per V/m \cite{bikson2004,deans2007} would scale the absolute values
by $0.24$--$0.36$ and $0.06$--$0.13$, respectively.

The detector compartment fixes the projection direction. A laminar direction for a pyramidal
axon initial segment, a local axon direction for a node, or an orientation
distribution for terminals defines a different access model. Isotropic averaging
is one possible distribution, not an anatomical consequence of choosing terminals.
Maximizing a geometric map over orientation gives an upper bound for that map
at each location; it does not identify the detector's orientation.

The projection direction $\mathbf n(\mathbf x)=\mathbf n^*(\mathbf x)$ is the envelope-maximizing direction at
the reference ratio 1:1: the closed-form maximizer, over unit vectors, of
$2\min(|\mathbf F_1\!\cdot\!\mathbf n|,|\mathbf F_2\!\cdot\!\mathbf n|)$
\cite{grossman2017}. It is chosen once at each node and then held fixed for 1:1, 1:3
and every ratio of the sweep, and all three maps use it. Along it, the 1:1 envelope map
equals its maximum over orientation. At any other ratio it is the envelope seen by an
element with that fixed orientation, which can be below that ratio's own maximum. As an
orientation sensitivity, every statistic was repeated along a second fixed field, the
radial direction from the brain centroid. The isotropic comparison averages the signed
mixing term over uniformly distributed orientations before its magnitude is taken:
$\langle\varepsilon_1\varepsilon_2\rangle_{\mathbf n}
=\lambda_0^2I_1^{\rm pair}I_2^{\rm pair}\,\mathbf F_1\!\cdot\!\mathbf F_2/3$, so the
isotropic beat excursion is
$\tfrac23\lambda_0^2I_1^{\rm pair}I_2^{\rm pair}|\mathbf F_1\!\cdot\!\mathbf F_2|$.
Elements whose beats have opposite phase cancel in this average.

The target is the left hippocampus of a FreeSurfer~6.0.0 \texttt{recon-all}
segmentation of the same T1 image (\texttt{aseg} label 17, 3903~mm$^3$)
\cite{fischl2002}. The region of interest (ROI) is the 1{,}568 gray- and white-matter
nodes inside it. The reference is the 273{,}544 gray-matter nodes outside the
ROI, including other deep gray matter. It is a gray-matter reference, not a cortex-only contrast.
The mesh and the segmentation share one image frame; tissue-volume, centroid and
white-matter checks confirmed their alignment. As a label sensitivity, the
Neuromorphometrics atlas hippocampus, warped into the subject with SimNIBS~3.2.3
\texttt{mni2subject} and the subject's nonlinear MNI transform (1{,}345 nodes; Dice 0.72
with the FreeSurfer label), replaced the ROI. As a tissue sensitivity, the ROI was
restricted to its gray-matter nodes (1{,}174 with the FreeSurfer label, 1{,}013 with
the atlas label).

The target-to-reference ratio is the ROI 95th percentile divided by the reference
95th percentile; its mean-based companion divides the ROI mean by the same
reference percentile. All percentiles and centroids weight nodes equally.
The quadratic excursion and $\mathcal I$ are second order in the field, whereas
the envelope excursion is first order. To separate this difference in degree
from the spatial weighting, each second-order statistic is also given
degree-matched, computed on the square root of the map. This is a geometric
comparison; the raw maps remain the respective detector predictions before gain.
The focus is the centroid of the top 10\% of ROI nodes,
ranked by the map's own value, with 5, 25 and 50\% as a sensitivity. Ranking makes the
focus exactly invariant to a uniform positive factor. Its displacement from 1:1 is
measured along the hippocampal long axis, the first principal axis of the ROI nodes,
with anterior positive. Field projections still use $\mathbf n^*$, independently
of this reporting axis. The 1:1 to 1:3 displacement was computed for every combination
of montage, label, ROI tissue definition and focus fraction (32 per map), each montage
along its own 1:1 direction.

No signed DC-source or mean-response map is computed. For a spatially uniform detector
coefficient, $Y_{\rm DC}=\tfrac12\Sig''(v^*)\mathcal I$ (\S\ref{sec:dc}) is
$\mathcal I$ rescaled, with the sign of $\Sig''(v^*)$. A nonuniform coefficient would
require an operating-point map, which a head model does not provide, and a mean
response would additionally require a circuit model. In particular,
$\Sig''(v^*)$ is curvature at the population model's operating point, not a
measurement of the transfer at the axon initial segment. Assigning that
coefficient to an anatomical detector requires a calibrated access and
input--output model.
